\section{Conclusiones}

Tras el procesamiento, análisis metrológico y discusión físico-matemática de las mediciones experimentales obtenidas en el ensayo de rendimiento del transformador monofásico de laboratorio, se establecen las siguientes conclusiones técnicas:

\begin{enumerate}
    \item \textbf{Eficiencia Elevada en Regímenes de Alta Carga:}
    Se constató experimentalmente que el transformador monofásico presenta un excelente rendimiento energético en la zona de alta potencia activa, alcanzando un valor de $\eta_{\text{med}} = 95{,}45 \pm 3{,}14\%$ a 1~kW y $\eta_{\text{med}} = 95{,}00 \pm 3{,}45\%$ a 800~W. En estos puntos de operación, la potencia útil suministrada a la carga supera con creces la suma de las pérdidas internas en el núcleo de hierro y en los conductores de cobre ($P_{\text{perd}} \approx 20\text{ W}$), situando a la máquina en la meseta óptima de transferencia de energía.

    \item \textbf{Límites Metrológicos del Método Directo en Bajas Cargas:}
    La obtención de rendimientos nominales de 100\% en las cargas reducidas de 600~W, 400~W y 200~W evidenció las limitaciones de sensibilidad de los vatímetros analógicos electrodinámicos con constante de escala $K_w = 10\text{ W/div}$ (obtenida de la placa del instrumento considerando el factor 5 con el TC $2:1$ en primario y el factor directo 10 en secundario). En bajas potencias, las pérdidas reales en vacío (15 a 25~W) quedan enmascaradas dentro de la incertidumbre instrumental de una sola división ($\pm 10\text{ W}$), provocando lecturas idénticas en primario y secundario. Este resultado justifica la necesidad metrológica de determinar el rendimiento por el método indirecto a partir de los parámetros del modelo circuital en lugar del método directo de carga cuando se opera a baja potencia.

    \item \textbf{Impacto Crítico del Factor de Potencia en Carga Inductiva:}
    El ensayo con carga inductiva ($L1-2$) demostró empíricamente la severa penalización que ocasiona un bajo factor de potencia en las máquinas eléctricas. A pesar de demandar la mayor corriente de toda la prueba ($I_1 = 6{,}0\text{ A}$, superando la capacidad nominal de placa de $5\text{ A}$), la potencia activa útil transferida colapsó a apenas $50\text{ W}$ debido a un factor de potencia secundario extremadamente bajo ($\cos\varphi_2 = 0{,}095$). La corriente circulante fue absorbida casi en su totalidad para excitar el campo magnético inductivo ($Q \approx 525\text{ VAR}$), sobrecargando térmicamente los conductores por efecto Joule sin aportar potencia activa útil, lo que ratifica la necesidad de corregir el factor de potencia en redes industriales.

    \item \textbf{Validación del Método de la Resistencia para Evaluación Térmica:}
    La aplicación de la ley de inferencia térmica normalizada para conductores de cobre, conforme al punto 5.7 de la norma COVENIN 3172-95 y ANSI/IEEE C57.12.91-2001, comprobó la idoneidad del método de la resistencia para estimar la temperatura media interna de los arrollamientos sin alterar la estructura del transformador. El devanado de baja tensión experimentó una elevación de temperatura de $\Delta T = 49{,}85 \pm 11{,}2~^{\circ}\mathrm{C}$ ($T_2 = 76{,}85~^{\circ}\mathrm{C}$), mientras que el de alta tensión registró $\Delta T = 30{,}89 \pm 14{,}6~^{\circ}\mathrm{C}$ ($T_2 = 57{,}89~^{\circ}\mathrm{C}$). La diferencia entre ambos responde al mayor confinamiento térmico y densidad de corriente del arrollamiento interno de BT. Ambas elevaciones se mantuvieron dentro de los límites admisibles para materiales aislantes Clase A ($60~^{\circ}\mathrm{C}$) fijados por COVENIN 536-94.

    \item \textbf{Segregación Analítica de Pérdidas y Validación del Modelo Circuital:}
    La exportación de los parámetros del modelo circuital determinados en la caracterización previa ($R_c = 240\,\Omega$, $R_{cc} = 1{,}25\,\Omega$) permitió discriminar cuantitativamente las pérdidas en el hierro ($P_{\text{Fe}} \approx 47 - 58\text{ W}$) y en el cobre ($P_{\text{Cu}} \le 5\text{ W}$ en la zona resistiva ensayada). Se constató que durante toda la prueba resistiva dominaron ampliamente las pérdidas constantes del núcleo ($P_{\text{Fe}} \gg P_{\text{Cu}} $), situando la operación en el tramo ascendente de la curva de eficiencia. Asimismo, el método indirecto demostró ser indispensable para revelar el colapso de rendimiento real a baja potencia ($45{,}88\%$ a 200~W) y bajo carga inductiva ($43{,}40\%$), subsanando las lecturas analógicas aparentemente ideales del método directo.
\end{enumerate}
